\documentclass{article}
\usepackage[utf8]{inputenc}
\usepackage{amsmath,amssymb}
\usepackage[most]{tcolorbox}
\usepackage{geometry}
\geometry{margin=2.5cm}

\newcommand{\step}[1]{\par\noindent\hangindent=3.5em\hangafter=1%
  \makebox[3.5em][l]{\textbf{#1.}}}
\newcommand{\steptag}[1]{\hfill{\small\textsf{[#1]}}}

\newtcolorbox{proofbox}[1]{
  colback=gray!5, colframe=gray!60,
  fonttitle=\bfseries, title={#1},
  breakable, enhanced,
  left=4pt, right=4pt, top=4pt, bottom=4pt,
  before skip=10pt, after skip=10pt
}

\begin{document}

\begin{proofbox}{There is a universal e\_H\textasciicircum{}r > 0 such that, for every finit...}
\step{1} There is a universal e\_H\textasciicircum{}r > 0 such that, for every finite-dimensional exact-unit epsilon\_r-C*-algebra with 0 <= epsilon\_r <= e\_H\textasciicircum{}r, the scalar-equivariant mu, sigma and the jointly continuous projected straight paths descend to breve-calU; the descended multiplication makes it a connected H-space, and the descended smooth map breve-sigma is a left inversion. \steptag{validated} \\[6pt]
\step{1.1} Universal ledger specialization and admissible operations/paths. Fix the universal witness tuple W from lem-stage1-polar-constant-ledger and put e\_pol\textasciicircum{}r := C\_rect*e\_S1 and e\_H\textasciicircum{}r := min\{e\_pol\textasciicircum{}r,e\_quot\textasciicircum{}r\}, where e\_quot\textasciicircum{}r is from lem-stage1-quotient-manifold-package. Then e\_H\textasciicircum{}r>0. Given 0<=epsilon\_r<=e\_H\textasciicircum{}r, set epsilon\_X:=epsilon\_r/C\_rect, delta:=delta\_*, q:=C\_grp*epsilon\_r, r\_-:=delta\_*-C\_pol*(epsilon\_r*delta\_*+delta\_*\textasciicircum{}2), and eta:=C\_path*(q+epsilon\_r*q+q\textasciicircum{}2). Since epsilon\_X<=e\_S1, clause (R) of lem-stage1-polar-constant-ledger gives all (A\_2),(A\_4),(A\_5),(A\_6) guards, in particular q<=1, C\_path*q<=1/4, q<r\_-, and eta<r\_-. Hence (A\_5) gives global C\textasciicircum{}1 maps mu(U,V)=u\_delta(U bold-dot V), sigma(U)=u\_delta(U\textasciicircum{}dagger), exact identities mu(J,U)=mu(U,J)=U and sigma(J)=J, and ||mu(sigma(U),U)-J||<=q; and (A\_6) gives, for every U\_0,U\_1 with ||U\_1-U\_0||<=q, the jointly continuous projected straight path H(t,U\_0,U\_1)=u\_delta((1-t)U\_0+tU\_1), with the stated endpoints and H(t,cU\_0,cU\_1)=cH(t,U\_0,U\_1). \steptag{validated} \\[4pt]
\step{1.2} Smoothness and scalar covariance of the same operations. At delta=delta\_*, (A\_2) of lem-stage1-polar-constant-ledger supplies the unique C\textasciicircum{}1 graph family covering calU and ||D\_\{A\textasciicircum{}perp\}f\_V-I||<1, so each normal derivative is invertible; lem-stage1-smooth-unitary-atlas upgrades exactly those graphs/charts to a smooth embedded atlas without changing points or first derivatives. Clause (A\_4) supplies the displayed C\textasciicircum{}1 polar diffeomorphism and inverse, and lem-stage1-smooth-polar-inverse upgrades those same maps to smooth maps. Thus every antecedent of lem-stage1-explicit-smooth-unitary-operations holds for the very u\_delta, mu and sigma fixed above; that lemma makes the same mu and sigma smooth and gives mu(cU,dV)=c*d*mu(U,V) and sigma(cU)=conj(c)*sigma(U) for all c,d in U(1). \steptag{validated} \\[4pt]
\step{1.3} Restriction to the identity component and descent. Let p:calU\_e -> breve-calU=calU\_e/U(1) be the orbit map and breve-e=[J]. The connected images mu(calU\_e x calU\_e) and sigma(calU\_e) contain J by the exact basepoint identities, hence lie in the connected component calU\_e; every admissible H-path whose initial point is in calU\_e also lies in calU\_e. The covariance identities therefore define continuous maps breve-mu([U],[V]):=[mu(U,V)] and breve-sigma([U]):=[sigma(U)], independent of representatives, and common-phase covariance makes each admissible projected straight path t $|-\rangle$ [H(t,U\_0,U\_1)] independent of a simultaneous replacement (U\_0,U\_1) by (cU\_0,cU\_1). For N>1 use the smooth quotient structure supplied by lem-stage1-quotient-manifold-package; for N=1, exact unitality and complex dimension one give calX=CJ, (cJ)\textasciicircum{}dagger=conj(c)J, (cJ) bold-dot (dJ)=cdJ, hence calU\_e=U(1)J and breve-calU is the one-point smooth quotient. Thus the stated descents exist in every finite dimension. \steptag{validated} \\[4pt]
\step{1.4} H-space law. The descended multiplication breve-mu is continuous, breve-mu(breve-e,breve-e)=breve-e, and the exact identities mu(J,U)=mu(U,J)=U imply breve-mu(breve-e,x)=x=breve-mu(x,breve-e). Hence both unit maps are literally the identity (and therefore basepoint-preserving homotopic to it by the constant homotopy), so (breve-calU,breve-e,breve-mu) is an H-space in the precise sense of def-h-space-left-inversion. \steptag{validated} \\[4pt]
\step{1.5} Connectedness. The space calU\_e is connected by its definition as the connected component of calU containing J, and p is continuous and surjective; therefore breve-calU=p(calU\_e) is connected (also consistently with lem-stage1-quotient-manifold-package when N>1 and with the singleton description when N=1). \steptag{validated} \\[4pt]
\step{1.6} Smoothness of the descended inversion. For N>1, equip breve-calU=calU\_e/U(1) with the canonical smooth orbit-quotient structure obtained directly from the smooth free scalar action on the smooth manifold calU\_e: the orbit map p is a smooth submersion with smooth local sections. Since sigma is smooth, preserves calU\_e, and satisfies sigma(cU)=conj(c)sigma(U), the representative-independent map breve-sigma([U])=[sigma(U)] is smooth. For N=1, breve-calU is a point and breve-sigma is its unique smooth self-map. \steptag{validated} \\[6pt]
\step{1.6.1} Canonical quotient charts from the allowed smooth action. By node 1.2, calU\_e is an open smooth submanifold of calU and alpha(c,U)=cU is smooth. The action is free: cU=U implies (c-1)U=0, while a unitary U has a right inverse and hence U is nonzero, so c=1. Fix U. The orbit tangent is the nonzero vector iU. Choose a linear complement E\_U of R(iU) in T\_U calU\_e and a smooth chart psi from a neighborhood of 0 in T\_U calU\_e to calU\_e with psi(0)=U and Dpsi\_0=id. For F(theta,v)=exp(i theta)psi(v), v in E\_U, one has DF\_(0,0)(a,w)=a iU+w, an isomorphism; the inverse-function theorem therefore gives a local product chart. Shrink its slice S\_U=psi(B cap E\_U) so that no phase outside the chosen identity arc carries S\_U into itself: otherwise shrinking neighborhoods would give sequences x\_n,y\_n->U and c\_n outside that arc with c\_n x\_n=y\_n; compactness of U(1) gives a subsequence c\_n->c outside the arc and then cU=U, contradicting freeness. Combining this separation with uniqueness in the local product chart shows that U(1) x S\_U -> U(1)S\_U is bijective modulo the evident orbit coordinate and is locally a diffeomorphism. Hence p(S\_U) is open, p|S\_U is a homeomorphism, and the charts (p|S\_U)\textasciicircum{}(-1) give the orbit set a smooth atlas: on an overlap, the unique phase carrying one slice representative to the other is smooth by either local product chart, so the transition map is smooth. In these charts p has local form (theta,v)$|-\rangle$v; consequently p is a smooth submersion and the inverse slice maps s\_U:p(S\_U)->S\_U are smooth local sections. This constructs the needed quotient smooth structure from allowed inputs, without asserting that the bare manifold-existence conclusion of lem-stage1-quotient-manifold-package already exported it. \steptag{validated} \\[4pt]
\step{1.6.2} Smooth equivariant descent. By nodes 1.2 and 1.3, sigma:calU\_e->calU\_e is smooth and sigma(cU)=conj(c)sigma(U), so breve-sigma([U]):=[sigma(U)] is well-defined. Let O=p(S\_U) be any quotient chart from the preceding construction and s\_U:O->calU\_e its smooth section. The identity breve-sigma|\_O = p o sigma o s\_U holds pointwise. Each factor is smooth (p by the quotient charts, sigma by node 1.2, and s\_U by construction), hence breve-sigma is smooth on every O and therefore globally. For N=1, node 1.3 identifies breve-calU with a singleton, whose unique self-map is smooth. \steptag{validated} \\[4pt]
\step{1.7} Left-inversion homotopy. For U in calU\_e set A(U):=mu(sigma(U),U). The covariance from lem-stage1-explicit-smooth-unitary-operations gives A(cU)=mu(conj(c)sigma(U),cU)=A(U), while (A\_5) of lem-stage1-polar-constant-ledger gives ||A(U)-J||<=q. Thus (A\_6) applies to (A(U),J), and P(t,[U]):=[H(t,A(U),J)] is a well-defined continuous homotopy from x $|-\rangle$ breve-mu(breve-sigma(x),x) to the constant breve-e. It is basepoint-preserving because A(J)=mu(sigma(J),J)=J and H(t,J,J)=J. Also breve-sigma(breve-e)=breve-e. Therefore breve-sigma is a left inversion map by def-h-space-left-inversion. \steptag{validated} \\[4pt]
\end{proofbox}

\end{document}
